\documentclass[11pt]{article}

\usepackage{lmodern}
\usepackage[T2A]{fontenc}
\usepackage[utf8]{inputenc}
\usepackage[russian]{babel}

\usepackage{amsmath, amssymb, mathtools}
\usepackage{geometry}
\geometry{a4paper, margin=2.2cm}
\usepackage{array}
\usepackage{booktabs}
\usepackage{longtable}
\usepackage{xcolor}
\usepackage{enumitem}
\usepackage{tcolorbox}
\tcbuselibrary{skins, breakable}

\definecolor{accent}{HTML}{1F4E79}
\definecolor{soft}{HTML}{EAF0F7}
\definecolor{warn}{HTML}{B5651D}
\definecolor{warnbg}{HTML}{FBF1E6}
\definecolor{okbg}{HTML}{ECF6EE}
\definecolor{okline}{HTML}{2E7D32}

\usepackage[unicode, colorlinks=true, linkcolor=accent, urlcolor=accent,
            citecolor=accent, bookmarksnumbered=true, bookmarksopen=true]{hyperref}

% ---- aesthetic boxes -------------------------------------------------------
\newtcolorbox{idea}[1][]{breakable, colback=soft, colframe=accent,
  boxrule=0.6pt, arc=2pt, left=8pt, right=8pt, top=6pt, bottom=6pt,
  fonttitle=\bfseries, title=#1}

\newtcolorbox{warnbox}[1][Замечание преподавателя]{breakable, colback=warnbg,
  colframe=warn, boxrule=0.6pt, arc=2pt, left=8pt, right=8pt, top=6pt, bottom=6pt,
  fonttitle=\bfseries, title=#1}

\newtcolorbox{okbox}[1][Вывод]{breakable, colback=okbg, colframe=okline,
  boxrule=0.6pt, arc=2pt, left=8pt, right=8pt, top=6pt, bottom=6pt,
  fonttitle=\bfseries, title=#1}

\newcommand{\F}{\mathbb{F}}
\renewcommand{\arraystretch}{1.25}

\usepackage{titlesec}
\titleformat{\section}{\Large\bfseries\color{accent}}{\thesection.}{0.6em}{}
\titleformat{\subsection}{\large\bfseries\color{accent}}{\thesubsection.}{0.5em}{}

\begin{document}

% ======================= TITLE ============================================
\begin{titlepage}
\centering
\vspace*{2.2cm}
{\Huge\bfseries\color{accent} Новые компьютерные технологии\par}
\vspace{0.8cm}
{\LARGE Таблицы Кэли, конечные поля\\[4pt] и сопряжённое пространство $V^{*}$\par}
\vspace{1.4cm}
{\large\itshape Учебно-пояснительная методичка:\\
теория, полное решение задач и разбор замечаний\par}
\vspace{2.4cm}
\begin{tcolorbox}[colback=soft, colframe=accent, width=0.8\textwidth, arc=3pt]
\centering
Документ отвечает на все вопросы преподавателя и разбирает\\
каждый пункт задания «с нуля» — чтобы материал стал понятен,\\
а не просто «посчитан кодом».
\end{tcolorbox}
\vfill
{\large 2026\par}
\end{titlepage}

% ======================= TOC ==============================================
\tableofcontents
\thispagestyle{empty}
\clearpage

% ##########################################################################
\section{Как читать этот документ}

Преподаватель прислал не «разнос кода», а \textbf{список вопросов к защите} и одну
главную мысль: код, который печатает правильные числа, ещё не показывает, что вы
\emph{понимаете} предмет. Поэтому документ построен так:

\begin{itemize}[leftmargin=1.4em]
  \item \textbf{Раздел 2} — прямой разбор каждого замечания преподавателя и готовые ответы.
  \item \textbf{Раздел 3} — Задача~I: что такое таблицы Кэли, как из них вырастают
        полугруппы, моноиды, группы, кольца и поля; все подсчёты для порядков $2,3,4$.
  \item \textbf{Раздел 4} — Задача~II: векторы, ковекторы, кодирование координатами,
        матрицы перехода, сопряжённое пространство $V^{*}$, изотропные векторы.
  \item \textbf{Раздел 5} — что конкретно поправить в коде и шпаргалка к защите.
\end{itemize}

Все числовые результаты — настоящие, полученные перебором в вашей программе
(\texttt{out/\allowbreak algebra\_\allowbreak counts.csv}). Цель — чтобы по каждому числу вы могли сказать
\emph{почему} оно такое.

% ##########################################################################
\section{Разбор замечаний преподавателя}

\begin{warnbox}[Что написал преподаватель (в сокращении)]
«Код, который что-то делает, не был самоцелью курса. Что такое векторы и ковекторы?
Как именно происходит кодирование? Что такое матрица перехода (и почему они у вас
описаны с точностью до наоборот)? Почему выбран вариант с 3-элементным полем и
размерностью три? Да ещё и с одним конкретным полем?»
\end{warnbox}

\subsection{«Код — не самоцель» (пункт 0)}

Курс — про \emph{смысл} объектов, а не про программу. Работающий перебор показывает
\emph{что} получилось ($6$ полей порядка~$3$ и т.\,д.), но не показывает, \emph{почему}
и \emph{что} эти объекты собой представляют. На защите ждут, что вы объясните понятия
своими словами. Ниже — ровно эти объяснения.

\subsection{Что такое вектор и что такое ковектор}

\begin{idea}[Определения]
\textbf{Вектор} — элемент векторного пространства $V$ над полем $\F$. Это
самостоятельный объект; у него \emph{нет} координат, пока не выбран базис.

\textbf{Ковектор} (линейный функционал, $1$-форма) — это \emph{линейное отображение}
$\varphi\colon V \to \F$. Множество всех таких отображений само образует векторное
пространство — \textbf{сопряжённое пространство} $V^{*}$.
\end{idea}

Ключевая мысль, которую проверяет преподаватель: вектор и ковектор живут в
\emph{разных} пространствах ($V$ и $V^{*}$) и при смене базиса ведут себя
\emph{по-разному} (вектор — контравариантно, ковектор — ковариантно, разделы
\ref{ssec:contra} и \ref{ssec:cova}).

\begin{warnbox}[Где это «провисает» в коде]
В программе вектор и ковектор — один и тот же тип \texttt{Vector} (\texttt{[]int}).
Машина не видит разницы между «точкой пространства $V$» и «функцией на $V$». Поэтому по
коду нельзя ответить на вопрос «что такое ковектор» — различие нужно вернуть хотя бы на
уровне понятий (см. раздел~\ref{sec:fix}).
\end{warnbox}

\subsection{Как происходит «кодирование»}

Это центральный вопрос. \emph{Кодирование} — это переход от абстрактного вектора к
набору чисел через выбранный базис.

\begin{idea}[Кодирование вектора координатами]
Пусть $e_1,\dots,e_n$ — базис $V$. Любой вектор $v\in V$ единственным образом
записывается как
\[
  v = x^1 e_1 + x^2 e_2 + \dots + x^n e_n .
\]
Набор $[v]_e=(x^1,\dots,x^n)\in\F^{n}$ называется \textbf{координатами} (кодом) вектора
$v$ в базисе $e$. Само отображение $v\mapsto [v]_e$ — это \emph{изоморфизм}
$V \xrightarrow{\sim} \F^{n}$.
\end{idea}

Главное: \textbf{код зависит от базиса, а сам вектор — нет}. Один вектор в разных
базисах имеет разные координаты. Именно это «двойное» представление (объект $\leftrightarrow$
его код) — суть второй части задания.

\begin{warnbox}[Где это «провисает» в коде]
Программа сразу порождает кортежи чисел (\texttt{AllVectors} строит наборы по основанию
$n$) — то есть работает уже \emph{в координатах}, причём в стандартном базисе.
Абстрактного вектора, который потом \emph{кодируется}, в модели нет, поэтому шаг
кодирования невидим. На защите стоит показать, что вы понимаете: число — это код, а не
сам вектор.
\end{warnbox}

\subsection{Матрица перехода и почему «наоборот»}
\label{ssec:transwarn}

Это самое конкретное замечание, и оно \emph{по делу}. Есть две разные матрицы, которые
легко перепутать:

\begin{enumerate}[leftmargin=1.6em]
  \item \textbf{Матрица перехода базисов} $C$ (от $e$ к $f$): её столбцы — это
        координаты \emph{новых} базисных векторов $f_j$ в \emph{старом} базисе $e$.
  \item \textbf{Матрица преобразования координат}: она равна $C^{-1}$ и переводит
        старые координаты вектора в новые.
\end{enumerate}

В коде и на странице матрицей «$A\to B$» названа как раз \emph{вторая} (преобразование
координат), а «$B\to A$» — \emph{первая}. По стандартному определению это \emph{поменяно
местами}. Подробный разбор на числах — в разделе~\ref{ssec:trans}; там видно, что
\emph{вычисления у вас верные и согласованные}, неверны только \emph{названия} матриц
относительно общепринятой договорённости. Именно это и значит «описаны с точностью до
наоборот».

\subsection{Почему $\F_3$ и размерность $3$ — и только одно поле}

Задание~II просит построить $V$ \emph{по заданному} полю $\F$ и размерности $\dim V\le 4$,
то есть метод должен работать для \emph{любого} поля из части~I и \emph{любой} размерности
$\le 4$. В отчёте же показан единственный случай $\F_3,\ \dim 3$.

\begin{idea}[Как закрыть вопрос]
Есть два честных пути:
\begin{itemize}[leftmargin=1.4em]
  \item \emph{Обобщить демонстрацию}: прогнать и показать примеры для $\F_2$ и $\F_4$ и
        для $\dim=2,4$ — код это уже умеет (флаги \texttt{-field}, \texttt{-dim}).
  \item \emph{Обосновать выбор}: $\F_3$, $\dim 3$ — минимальный нетривиальный пример,
        где видно всё сразу: ненулевые недиагональные координаты при смене базиса,
        нетривиальный дуальный базис и непустое множество изотропных векторов
        (в $\F_2$ многое вырождается). Это законный аргумент, если его проговорить.
\end{itemize}
Важно: над разными полями результаты различаются (число изотропных векторов, наличие
квадратных корней из $-1$ и т.\,п.), поэтому «одно конкретное поле» само по себе не доказывает
общий случай.
\end{idea}

% ##########################################################################
\section{Задача I. Таблицы Кэли и конечные поля}

\subsection{Бинарная операция и таблица Кэли}

Зафиксируем конечное множество $S=\{0,1,\dots,n-1\}$. \textbf{Бинарная операция} — это
функция $*\colon S\times S\to S$. Её удобно задать \textbf{таблицей Кэли}: в клетке на
пересечении строки $a$ и столбца $b$ стоит значение $a*b$.

\begin{idea}[Сколько всего таблиц для одной операции — пункт 1а]
В таблице $n\times n$ клеток, и в каждую можно независимо поставить любой из $n$
элементов. Значит всего операций (таблиц)
\[
  N(n)=n^{\,n^{2}}.
\]
\[
  N(2)=2^{4}=16,\qquad N(3)=3^{9}=19683,\qquad N(4)=4^{16}=4\,294\,967\,296.
\]
\end{idea}

Дальше мы среди всех этих таблиц отбираем те, что удовлетворяют дополнительным аксиомам.

\subsection{Иерархия структур с одной операцией (пункты 1б–1д)}

\begin{description}[leftmargin=2.2em, style=nextline]
  \item[Полугруппа] операция \emph{ассоциативна}: $(a*b)*c=a*(b*c)$ для всех $a,b,c$.
  \item[Моноид] полугруппа, в которой есть \emph{нейтральный} элемент $e$:
        $e*a=a*e=a$.
  \item[Группа] моноид, в котором у каждого $a$ есть \emph{обратный} $a^{-1}$:
        $a*a^{-1}=a^{-1}*a=e$.
  \item[Абелева группа] группа с \emph{коммутативной} операцией: $a*b=b*a$.
\end{description}

Эти классы вложены: $\text{группы}\subset\text{моноиды}\subset\text{полугруппы}$. Поэтому
«полугрупп, не являющихся моноидами» $=$ (полугруппы) $-$ (моноиды), и аналогично дальше.

\begin{idea}[Результаты перебора (одна операция)]
\centering
\begin{tabular}{l r r r}
\toprule
\textbf{Класс} & $n=2$ & $n=3$ & $n=4$\\
\midrule
Всего таблиц $n^{n^2}$        & 16 & 19683 & 4\,294\,967\,296\\
Полугрупп                     & 8  & 113   & 3\,492\\
\quad из них не моноидов      & 4  & 80    & 2\,868\\
Моноидов                      & 4  & 33    & 624\\
\quad из них не групп         & 2  & 30    & 608\\
Групп                         & 2  & 3     & 16\\
\quad из них неабелевых       & 0  & 0     & 0\\
Абелевых групп                & 2  & 3     & 16\\
\bottomrule
\end{tabular}
\end{idea}

\noindent\textbf{Как это читать.} Для $n\le 4$ \emph{все} группы абелевы (наименьшая
неабелева группа $S_3$ имеет порядок~$6$). Числа групп $2,3,16$ — это количество
\emph{размеченных} таблиц (на фиксированном множестве с конкретными метками), а не классов
изоморфизма: например, порядок~$4$ даёт $2$ группы с точностью до изоморфизма
($\mathbb{Z}_4$ и $\mathbb{Z}_2\times\mathbb{Z}_2$), но $16$ размеченных таблиц.

\subsection{Две операции: от полуколец до полей (пункты 2а–2д)}

Теперь на множестве заданы \emph{две} операции — сложение $+$ и умножение $\cdot$.

\begin{description}[leftmargin=2.2em, style=nextline]
  \item[Полукольцо] $(S,+)$ — коммутативный моноид с нулём $0$; $(S,\cdot)$ —
        полугруппа; выполнены обе дистрибутивности
        $a(b+c)=ab+ac$, $(a+b)c=ac+bc$; ноль поглощает: $0\cdot a=a\cdot 0=0$.
  \item[Кольцо] полукольцо, в котором $(S,+)$ — \emph{абелева группа}
        (т.\,е. у сложения есть вычитание).
  \item[Кольцо с единицей] кольцо, где у умножения есть нейтральный элемент $1$.
  \item[Тело] кольцо с единицей ($1\neq 0$), в котором все ненулевые элементы образуют
        \emph{группу по умножению} (есть деление).
  \item[Поле] \emph{коммутативное} тело.
\end{description}

Вложенность: $\text{поля}\subset\text{тела}\subset\text{кольца с 1}\subset\text{кольца}
\subset\text{полукольца}$.

\begin{idea}[Результаты перебора (две операции)]
\centering
\begin{tabular}{l r r r}
\toprule
\textbf{Класс} & $n=2$ & $n=3$ & $n=4$\\
\midrule
Полуколец                       & 8 & 129 & 6\,508\\
\quad не являющихся кольцами     & 4 & 120 & 6\,348\\
Колец (всего)                   & 4 & 9   & 160\\
\quad колец без единицы          & 2 & 3   & 88\\
Колец с единицей                & 2 & 6   & 72\\
\quad не являющихся телами        & 0 & 0   & 60\\
Тел                             & 2 & 6   & 12\\
\quad не являющихся полями        & 0 & 0   & 0\\
Полей                           & 2 & 6   & 12\\
\bottomrule
\end{tabular}
\end{idea}

\noindent\textbf{Два важных наблюдения.}
\begin{itemize}[leftmargin=1.4em]
  \item \emph{«Тел, не являющихся полями» — всегда $0$.} Это теорема Веддербёрна:
        \emph{любое конечное тело коммутативно}, значит является полем. Для $n\le 4$ это
        видно прямо из перебора.
  \item \emph{«Колец с единицей, не являющихся телами» при $n=4$ равно $60$.} Пример —
        $\mathbb{Z}_4$ (кольцо вычетов): элемент $2$ не обратим, поэтому это кольцо с
        единицей, но не тело. При $n=2,3$ ($n$ простое) каждое кольцо с единицей без
        делителей нуля оказывается полем, поэтому там $0$.
\end{itemize}

\subsection{Конкретные поля порядков 2, 3, 4}

Поле порядка $n$ существует тогда и только тогда, когда $n=p^{k}$ — степень простого.
Числа $2,3,4=2^2$ подходят. Ниже — канонические представители (метки выбраны так, что
$0$ — нуль, $1$ — единица).

\paragraph{Поле $\F_2=\mathbb{Z}/2$.}
\[
\begin{array}{c|cc}
+ & 0 & 1\\\hline
0 & 0 & 1\\
1 & 1 & 0
\end{array}
\qquad\qquad
\begin{array}{c|cc}
\cdot & 0 & 1\\\hline
0 & 0 & 0\\
1 & 0 & 1
\end{array}
\]

\paragraph{Поле $\F_3=\mathbb{Z}/3$.}
\[
\begin{array}{c|ccc}
+ & 0 & 1 & 2\\\hline
0 & 0 & 1 & 2\\
1 & 1 & 2 & 0\\
2 & 2 & 0 & 1
\end{array}
\qquad\qquad
\begin{array}{c|ccc}
\cdot & 0 & 1 & 2\\\hline
0 & 0 & 0 & 0\\
1 & 0 & 1 & 2\\
2 & 0 & 2 & 1
\end{array}
\]

\paragraph{Поле $\F_4=\mathrm{GF}(4)$.}
Это \emph{не} $\mathbb{Z}/4$! Здесь характеристика $2$ (т.\,е. $x+x=0$), а элементы
$2,3$ — корни многочлена $x^2+x+1$. Обозначив $2=\alpha$, имеем $3=\alpha^2=\alpha+1$.
\[
\begin{array}{c|cccc}
+ & 0 & 1 & 2 & 3\\\hline
0 & 0 & 1 & 2 & 3\\
1 & 1 & 0 & 3 & 2\\
2 & 2 & 3 & 0 & 1\\
3 & 3 & 2 & 1 & 0
\end{array}
\qquad\qquad
\begin{array}{c|cccc}
\cdot & 0 & 1 & 2 & 3\\\hline
0 & 0 & 0 & 0 & 0\\
1 & 0 & 1 & 2 & 3\\
2 & 0 & 2 & 3 & 1\\
3 & 0 & 3 & 1 & 2
\end{array}
\]
Сложение — это «исключающее ИЛИ» по битам; умножение — степени образующего $\alpha=2$:
$\alpha^1=2,\ \alpha^2=3,\ \alpha^3=1$.

\subsection{Изоморфизм полей и формула \texorpdfstring{$n!/|\mathrm{Aut}|$}{n!/|Aut|} (пункт 2е)}

\begin{idea}[Изоморфизм]
Два поля $\F,\F'$ \textbf{изоморфны}, если есть биекция $\phi\colon\F\to\F'$, сохраняющая
обе операции: $\phi(a+b)=\phi(a)+\phi(b)$ и $\phi(ab)=\phi(a)\phi(b)$. По сути это «то же
поле, только с переименованными элементами».
\end{idea}

\textbf{Теорема.} Конечное поле заданного порядка \emph{единственно с точностью до
изоморфизма}. Поэтому все найденные поля порядка $n$ лежат в одном классе изоморфизма —
в таблице это столбец «поля изоморфны — да, классов $1$».

Откуда тогда $2,6,12$ \emph{размеченных} полей? Каждая перестановка меток множества
$\{0,\dots,n-1\}$ переносит таблицу в (возможно) другую таблицу того же поля. Перестановок
ровно $n!$. Часть из них переводит поле \emph{в себя} — это \textbf{автоморфизмы}; их число
$|\mathrm{Aut}(\F)|$. По теореме об орбите и стабилизаторе число различных размеченных
копий равно
\[
  \#\{\text{размеченных полей порядка }n\}=\frac{n!}{|\mathrm{Aut}(\F)|}.
\]

\begin{idea}[Сверка с перебором]
\centering
\begin{tabular}{c c c c c}
\toprule
$n$ & $n!$ & $|\mathrm{Aut}(\F_n)|$ & $n!/|\mathrm{Aut}|$ & найдено полей\\
\midrule
2 & 2  & 1 & 2  & 2\\
3 & 6  & 1 & 6  & 6\\
4 & 24 & 2 & 12 & 12\\
\bottomrule
\end{tabular}
\end{idea}

\noindent У $\F_2,\F_3$ автоморфизм только тождественный. У $\F_4$ их два: тождественный и
\emph{автоморфизм Фробениуса} $x\mapsto x^{2}$ (он переставляет $2\leftrightarrow 3$ и
оставляет $0,1$ на месте). Отсюда $24/2=12$.

% ##########################################################################
\section{Задача II. Пространство $V$ и сопряжённое $V^{*}$}

Везде ниже фиксируем поле $\F$ (одно из $\F_2,\F_3,\F_4$) и берём
$V=\F^{n}$, $n=\dim V\le 4$. Все примеры с числами — для $\F_3,\ n=3$ (этот случай и
разобран в коде), но определения и формулы \emph{не зависят} от выбора поля и размерности.

\subsection{Векторное пространство и базис (пункт 1а)}

$V$ — множество с операциями «$+$» и умножением на скаляр из $\F$, удовлетворяющими
обычным аксиомам. Набор $e_1,\dots,e_n$ — \textbf{базис}, если он
\emph{линейно независим} и \emph{порождает} $V$. Эквивалентно: матрица, составленная из
этих векторов как столбцов, имеет \emph{полный ранг} $n$ (обратима над $\F$). Именно так
проверяет код (метод Гаусса над $\F$, без \texttt{float}).

\textbf{Два базиса примера} ($\F_3,\ n=3$):
\[
A=\{e_1,e_2,e_3\}=\Big\{(1,0,0),(0,1,0),(0,0,1)\Big\},\quad
B=\{f_1,f_2,f_3\}=\Big\{(1,1,0),(0,1,1),(0,0,1)\Big\}.
\]
$A$ — стандартный; $B$ имеет полный ранг (определитель $=1\neq 0$ в $\F_3$), значит тоже
базис.

\subsection{Разложение вектора по базису (пункт 1б)}

Координаты $[v]_e$ находятся решением системы $\;[e_1\,|\,\dots\,|\,e_n]\,[v]_e=v\;$, то
есть умножением на обратную матрицу базиса. Для $v=(1,2,0)$ в $\F_3$:
\[
[v]_A=(1,2,0)\quad(\text{в стандартном базисе координаты совпадают с самим вектором}),
\]
\[
[v]_B=(1,1,2)\quad(\text{потому что } 1\!\cdot\!f_1+1\!\cdot\!f_2+2\!\cdot\!f_3=(1,2,0)\ \text{в }\F_3).
\]
Проверка последнего: $f_1+f_2+2f_3=(1,1,0)+(0,1,1)+2(0,0,1)=(1,2,3)=(1,2,0)$ по модулю $3$.

\subsection{Матрицы перехода (пункт 1в)}
\label{ssec:trans}

\begin{idea}[Аккуратное определение]
\textbf{Матрица перехода} $C$ от базиса $A$ к базису $B$ задаётся разложением
\emph{новых} векторов по \emph{старым}:
\[
  f_j=\sum_{i=1}^{n} C_{ij}\,e_i,
\]
то есть \textbf{столбцы $C$ — это координаты новых базисных векторов $f_j$ в старом
базисе $A$}. Тогда координаты любого вектора пересчитываются \emph{обратной} матрицей:
\[
  \boxed{\,[v]_A=C\,[v]_B\,}\qquad\Longleftrightarrow\qquad [v]_B=C^{-1}[v]_A .
\]
\end{idea}

Для нашего примера старый базис $A$ стандартный, поэтому столбцы $C$ — это просто векторы
$B$:
\[
C_{A\to B}=\big[\,f_1\ f_2\ f_3\,\big]=
\begin{pmatrix}1&0&0\\1&1&0\\0&1&1\end{pmatrix},
\qquad
C_{A\to B}^{-1}=C_{B\to A}=
\begin{pmatrix}1&0&0\\2&1&0\\1&2&1\end{pmatrix}\ (\text{в }\F_3).
\]
Проверка обратности: $C_{A\to B}\,C_{B\to A}=I$ (взаимно обратны — да).

\begin{warnbox}[Почему «наоборот» — на этих самых числах]
В коде/на странице под именем «\,$T\;(A\to B)$\,» стоит матрица
$\big(\begin{smallmatrix}1&0&0\\2&1&0\\1&2&1\end{smallmatrix}\big)$ — но это $C^{-1}$,
т.\,е. матрица \emph{преобразования координат} $A\!\to\!B$, она же \emph{стандартная
матрица перехода базисов} $B\to A$. А под «\,$T\;(B\to A)$\,» стоит
$\big(\begin{smallmatrix}1&0&0\\1&1&0\\0&1&1\end{smallmatrix}\big)=C_{A\to B}$.
То есть ярлыки переставлены. \textbf{Сами вычисления верны и согласованы} — поэтому все
проверки проходят, — но названия противоречат общепринятому определению «матрицы
перехода базисов». Достаточно поменять подписи (или явно оговорить свою договорённость).
\end{warnbox}

\subsection{Сопряжённое пространство $V^{*}$ и дуальный базис (пункт 2а)}

$V^{*}$ — пространство всех линейных функций $\varphi\colon V\to\F$. Если в $V$ выбран
базис $e_1,\dots,e_n$, то в $V^{*}$ есть \textbf{сопряжённый (дуальный) базис}
$e^1,\dots,e^n$, определённый правилом
\[
  e^{i}(e_j)=\delta^{i}_{j}=\begin{cases}1,& i=j,\\ 0,& i\neq j.\end{cases}
\]
То есть $e^{i}$ «считывает $i$-ю координату» вектора. В координатах: если базис $A$ записан
матрицей-столбцами $M=[e_1|\dots|e_n]$, то строки $M^{-1}$ — это координаты дуального
базиса.

Для примера ($B$): дуальный базис $B^{*}=\{f^1,f^2,f^3\}$ имеет координаты — строки
$C_{A\to B}^{-1}$:
\[
f^1=(1,0,0),\quad f^2=(2,1,0),\quad f^3=(1,2,1).
\]
Проверка сопряжённости: $f^{i}(f_j)=\delta^{i}_j$ — выполнено для всех $i,j$.

\subsection{Контравариантность координат вектора (пункт 2б)}
\label{ssec:contra}

При смене базиса $A\to B$ координаты вектора умножаются на $C^{-1}$:
\[
  [v]_B=C^{-1}[v]_A.
\]
Координаты меняются \emph{обратно} к базисным векторам — поэтому их называют
\textbf{контравариантными} (верхний индекс $x^{i}$). Проверка на числах:
$C^{-1}[v]_A=\big(\begin{smallmatrix}1&0&0\\2&1&0\\1&2&1\end{smallmatrix}\big)(1,2,0)^{\!\top}
=(1,1,2)^{\top}=[v]_B$. Сходится.

\subsection{Ковариантность координат ковектора (пункт 2в)}
\label{ssec:cova}

Для ковектора $\varphi$ с координатами $[\varphi]_{A^{*}}$ (в дуальном базисе $A^{*}$)
смена базиса работает \emph{прямой} матрицей, причём транспонированной:
\[
  [\varphi]_{B^{*}}=C^{\top}\,[\varphi]_{A^{*}}.
\]
Компоненты меняются \emph{так же}, как базисные векторы $V$ — поэтому их называют
\textbf{ковариантными} (нижний индекс $\varphi_i$). Возьмём ковектор $w$ с координатами
$(1,1,1)$ в $A^{*}$:
\[
[w]_{B^{*}}=C^{\top}(1,1,1)^{\top}=
\begin{pmatrix}1&1&0\\0&1&1\\0&0&1\end{pmatrix}(1,1,1)^{\top}=(2,2,1)^{\top}.
\]
Сходится с расчётом программы.

\begin{okbox}[Главный смысл части II]
Один и тот же «массив чисел» ведёт себя по-разному в зависимости от того, \emph{чем} он
является. Координаты вектора преобразуются через $C^{-1}$ (контравариантно), координаты
ковектора — через $C^{\top}$ (ковариантно). Различие вектор/ковектор — не формальность,
а именно \emph{закон преобразования}. Это и есть ответ на вопрос «что такое векторы и
ковекторы и как происходит кодирование».
\end{okbox}

\subsection{Изотропные векторы (пункт 2г)}

Зададим на $V$ симметричную билинейную форму матрицей Грама $G$. Вектор $v\neq 0$
называется \textbf{изотропным}, если $\langle v,v\rangle = v^{\top}G\,v = 0$.

Для единичной формы $G=I$ в $\F_3,\ n=3$ условие — $x_1^2+x_2^2+x_3^2\equiv 0 \pmod 3$.
Перебор ненулевых $v$ даёт $8$ изотропных векторов:
\[
(1,1,1),(2,1,1),(1,2,1),(2,2,1),(1,1,2),(2,1,2),(1,2,2),(2,2,2).
\]
\textbf{Зачем это в задании.} Над $\mathbb{R}$ из $x_1^2+\dots=0$ следует $v=0$ — изотропных
нет. Над конечным полем это уже не так: квадраты ведут себя иначе (в $\F_3$ имеем
$1^2=1,\ 2^2=1$), и ненулевые «нулевой длины» векторы появляются. Это наглядно показывает,
\emph{почему} выбор поля важен — и косвенно отвечает на вопрос преподавателя «почему именно
это поле»: над $\F_2$ или над $\mathbb{R}$ картина была бы другой.

% ##########################################################################
\section{Что поправить в коде и шпаргалка к защите}
\label{sec:fix}

\subsection{Минимальные правки в коде}
\begin{enumerate}[leftmargin=1.6em]
  \item \textbf{Переименовать матрицы перехода.} Привести подписи к стандарту: то, что
        сейчас называется $A\to B$, по определению есть матрица $B\to A$ (и наоборот).
        Либо оставить как есть, но \emph{явно} написать свою договорённость:
        «у нас матрица $A\to B$ — это матрица преобразования координат».
  \item \textbf{Развести вектор и ковектор.} Хотя бы ввести отдельный тип/обёртку
        \texttt{Covector} и комментарий, что это элемент $V^{*}=\mathrm{Hom}(V,\F)$, а
        применение к вектору — это $\varphi(v)=\sum \varphi_i v^{i}$.
  \item \textbf{Показать кодирование явно.} Вывести в отчёте шаг $v\mapsto[v]_e$ для
        \emph{произвольного} (заданного на входе) базиса, а не только для стандартного.
  \item \textbf{Снять привязку к $\F_3,\dim 3$.} Прогнать демонстрацию для $\F_2,\F_4$ и
        $\dim\in\{2,4\}$ и приложить результаты — флаги уже есть.
\end{enumerate}

\subsection{Шпаргалка: что спросят и что отвечать}
\begin{idea}[Семь ответов]
\begin{enumerate}[leftmargin=1.6em]
  \item \textbf{Вектор} — элемент $V$; \textbf{ковектор} — линейная функция $V\to\F$,
        элемент $V^{*}$.
  \item \textbf{Кодирование} — координаты $v=\sum x^i e_i$; код зависит от базиса,
        вектор — нет; $v\mapsto[v]_e$ — изоморфизм $V\cong\F^n$.
  \item \textbf{Матрица перехода} $C$ ($A\to B$): столбцы — новые базисные векторы в
        старом базисе; координаты пересчитываются как $[v]_B=C^{-1}[v]_A$.
  \item \textbf{Контравариантность}: координаты вектора идут через $C^{-1}$.
        \textbf{Ковариантность}: координаты ковектора идут через $C^{\top}$.
  \item \textbf{Дуальный базис}: $e^i(e_j)=\delta^i_j$; координаты — строки $M^{-1}$.
  \item \textbf{Поле единственно} с точностью до изоморфизма; число размеченных копий
        $=n!/|\mathrm{Aut}|$; у $\F_4$ автоморфизм Фробениуса $x\mapsto x^2$.
  \item \textbf{Теорема Веддербёрна}: конечное тело $=$ поле, поэтому «тел не полей» $=0$.
  \end{enumerate}
\end{idea}

\vspace{1em}
\begin{center}\itshape Конец методички.\end{center}

\end{document}
